\subsection{Full restriction orbits and fourth rotations for $R9\_2\_137$}
\label{sec:r137-restrictions}
The remaining ring has source dimensions $(1^4,2,d^4)$, where
$d=4+3\sqrt2$ and $D=144+96\sqrt2$. Work in the positive representative
supplied by the general-complex existence reduction. The four large
second indicators are $+1$ by the preceding Galois-pairing argument.

\begin{lemma}[Galois invariance of full restriction columns]
Let $\mathcal C$ be pseudo-unitary with commutative Grothendieck ring,
and let $K$ contain its complete character table and
$D=\FPdim(\mathcal C)$. For a simple $Z$ of $\mathcal Z(\mathcal C)$,
write $v_Z$ for the coefficient vector of $F(Z)$. Every modular Galois
automorphism fixing $K$ preserves $v_Z$.
If the square-Galois orbit of $\theta_Z$ under these automorphisms has
size $k$, the induction Gram matrix $G$ satisfies
\[
 k v_Zv_Z^{\mathsf T}\le G.
\]
When $G$ is positive definite, this is equivalent to
$k v_Z^{\mathsf T}G^{-1}v_Z\le1$.
\end{lemma}
\begin{proof}
For each irreducible character $\chi_E$ let $A_E$ be the corresponding
distinguished simple in $I(1)$, of dimension $D/f_E$. The center character
identity~\cite[Proposition 2.8 and Corollary 4.5]{Deshpande} and
\cite[Theorem 2.13]{Ostrik} gives, with normalized center matrix $s=S/D$,
\[
 s(A_E,Z)=\chi_E(FZ)/f_E\in K.
\]
The signed Galois permutation of columns obeys
$\sigma(s(A_E,Z))=\epsilon_\sigma(Z)s(A_E,\pi_\sigma Z)$.
The unit row is $\dim(Z)/D\in K$ and both it and its permuted entry are
positive, so $\epsilon_\sigma(Z)=1$. Every $\chi_E$ therefore takes the
same value on $FZ$ and $F(\pi_\sigma Z)$. The complete character table
separates the complexified commutative based algebra, proving equality
of their full coefficient vectors.

The center has canonical anomaly-one twist normalization, and modular
Galois symmetry gives $\theta_{\pi_\sigma Z}=\theta_Z^{a^2}$ when
$\sigma(\zeta)=\zeta^a$~\cite[Lemma 4.2]{DLN}. A $k$-element twist orbit
therefore provides at least $k$ distinct simples with column $v_Z$.
Subtracting their contributions from $G=\sum_Wv_Wv_W^{\mathsf T}$
proves the matrix inequality. Conjugation by $G^{-1/2}$ gives the final
equivalence.
\end{proof}

The nine distinct source characters of $R9\_2\_137$ all lie in
$K=\Q(\sqrt2)$; their multiplication identities and completeness are
checked against the entire tensor. Fixing only an FP-dimension field
would not justify the lemma in general. In particular a twist of order
divisible by nine supplies at least three copies of its entire
restriction column, not merely three objects of equal dimension.

The Galois-unit character formula also fixes a central simple $B$ with
\[
 F(B)=1+X_3+X_4+X_5+X_6,\qquad \dim B=17+12\sqrt2.
\]
The full character calculation identifies this vector: its values are
$f_E/\sigma(f_E)$ for $\sigma(\sqrt2)=-\sqrt2$.
Thus, for a high column with $v_Z(0)=0$, the stronger necessary bound is
\begin{equation}\tag{G137}\label{eq:r137-residual-gram}
 k v_Zv_Z^{\mathsf T}\le
 H:=G-e_0e_0^{\mathsf T}-v_Bv_B^{\mathsf T}.
\end{equation}
Here $H$ is positive definite, with determinant $6849298432$. The
distinguished low column for $B$ is imposed separately.

\begin{proposition}[Fourth-rotation bound]
For each large simple $X$ of $R9\_2\_137$,
\[
 \operatorname{Re}\nu_4(X),\ \operatorname{Im}\nu_4(X)
 \in\{-4,-2,0,2,4\}.
\]
This holds for both TY signs and does not require unitarity.
\end{proposition}
\begin{proof}
The source-local stabilizer anticommutator used in
\eqref{eq:source-cubic-clifford} gives four pair-insertion involutions
on $V=\operatorname{Hom}(1,X^4)$. Adjacent insertions anticommute,
opposite ones commute, and pivotal rotation $E$ cyclically permutes
them with $E^4=1$. The central products of opposite insertions split
the insertion algebra into four sign blocks. The two rotation-fixed
blocks give, respectively, opposite and adjacent fourth-root pairs;
the other two blocks are exchanged and give the regular four-root
spectrum, which is a sum of two adjacent pairs. This is an algebraic
Clifford decomposition and uses no inner product.

The tensor gives
\[
 \dim V=18,\qquad X^2=1+h+2\sum_{\text{large }Y}Y.
\]
All displayed second indicators are $+1$. Grouping four strands into
two pairs therefore yields
$\operatorname{tr}(E^2)=\nu_2(X^2)=10$.
Let $a_+,a_-$ count opposite pairs whose squared eigenvalue is
$+1,-1$, and let $b$ count adjacent pairs. Then
\[
 a_++a_-+b=9,\qquad 2(a_+-a_-)=10,\qquad b=4-2a_-\le4.
\]
Only adjacent pairs contribute to $\operatorname{tr}(E)$; each has
real and imaginary parts $\pm1$. The possible values are therefore
the five even integers displayed.
\end{proof}

\begin{proposition}[The center has no nontrivial tensor factor]
For a categorification of $R9\_2\_137$, its center is either simple,
or has fusion-subcategory lattice
\[
 \Vec,\quad \mathrm{sVec},\quad f',\quad \mathcal Z(\mathcal C),
\]
where $f$ is the central fermion and $f'$ is nonsplit and s-simple.
In particular its center is not a nontrivial Deligne product.
\end{proposition}
\begin{proof}
Every nonintegral source simple generates the full source, and all
proper source subcategories lie in the integral TY subcategory.
The integral-center theorem above says
$\mathcal Z(\mathcal C)_{\rm int}=\Vec$ or $\mathrm{sVec}$.
Every nonintegral central subcategory therefore forgets dominantly
to $\mathcal C$. Its dimension is strictly greater than $D$:
equality would make the tensor functor an equivalence and supply a
braiding on $\mathcal C$, whereas a braiding would lift its four
invertibles to a center having at most two.

A subcategory and its centralizer cannot both be nonintegral,
because their dimension product is $D^2$ by the modular centralizer
formula~\cite{DGNO}. If the integral part is $\Vec$, double
centralization proves simplicity. If it is $\mathrm{sVec}$,
the centralizer of any nonintegral proper subcategory must equal
$\mathrm{sVec}$, and double centralization identifies that
subcategory with $f'$. A splitting $f'=\mathcal M\boxtimes\mathrm{sVec}$
would produce a further proper nonintegral subcategory $\mathcal M$,
contradicting this list.
\end{proof}


\subsection{Full Galois complement and integral Verlinde restrictions}
\label{sec:r137-integral-verlinde}
Use source indices $0,\ldots,8$, with invertibles $0,1,2,7$, the
dimension-two object $8$, and $d=4+3\sqrt2$ at $3,4,5,6$.
For $v=[FZ]$ write $a=v_0+v_1+v_2+v_7+2v_8$,
$b=v_3+v_4+v_5+v_6$ and $\dim Z=a+bd$.
All nine source characters lie in $K=\mathbb Q(\sqrt2)$.
For its distinguished $A_E\subset I(1)$,
$s(A_E,Z)=\chi_E(v)/f_E$ in normalized modular notation.
Signed Galois symmetry and the exact character table give
\[
 (Tv)_i=v_i\ (i=0,1,2,7,8),\qquad
 (Tv)_i=a-v_i\ (i=3,4,5,6).
\]
For a Galois automorphism nontrivial on $K$,
$v_{\sigma Z}=\varepsilon_Z Tv$. If $a>0$, an unchanged positive
coordinate forces $\varepsilon_Z=1$, hence $v_i\le a$ for each large
coordinate. If $a=0$, positivity forces $\varepsilon_Z=-1$ and the
physical vector is unchanged. The squared-twist Galois rule preserves
all normalized relative trace types, since $\sigma^2$ fixes $K$.
Aggregate multiplicities of $(v,\text{trace type})$ and
$(Tv,\text{same type})$ therefore agree. Duplicate variables are grouped;
individual relaxation variables need not be matched. Missing partners
can be deleted because the inherited domain is complete.
The character table, all source products and $T^2=1$ are checked exactly.

Let $B$ be the Galois conjugate of the central unit and $D'=144-96\sqrt2$.
The distinguished-row identities force
\[
F(B)=1+X_3+X_4+X_5+X_6,\qquad
s(B,Z)=\frac{a+(4a-b)d}{D}.
\]
Verlinde gives nonnegative integers $n_E=N_{BB}^{A_E}$, namely
\[
n_E=\sum_Z
\frac{(a+(4a-b)d)^2\chi_E(v_Z)}{D(a+bd)f_E}.
\]
Their irrational coefficients vanish. In particular
$n_B=\sum_Z(a+(4a-b)d)^3/(D^2(a+bd))$.
These are linear restrictions on actual central-simple multiplicities.

\begin{theorem}\label{thm:r137-positive-ty-exclusion}
The positive Tambara--Yamagami sign is impossible in a positive
categorification of $R9\_2\_137$. The negative sign remains unresolved.
\end{theorem}
\begin{proof}
The complete cubic reduction above forces all four projected third
indicators to $1/2$ for the positive sign. No fourth indicator is fixed
to $2$: the four equations choosing that test value are removed.
On the resulting necessary system, full-column Galois complement and
the nine vanishing irrational parts, nine exact rational certificates
give the following successive bounds:
\[
\begin{array}{c|c|c}
\text{coefficient}&\text{bound}&\text{integer consequence}\\ \hline
n_B&1/2\le n_B\le31/18&n_B=1\\
\text{two codegree-16 }n_E&0\le n_E\le17/30&n_E=0\\
\text{four codegree-8 }n_E&0\le n_E\le3/4&n_E=0\\
\text{codegree-4 }n_E&0\le n_E\le1&n_E\le1 .
\end{array}
\]
Each later stage uses the earlier proved integral equalities.
There is no arbitrary upper bound on integral variables. For every
certificate, its equality multipliers, nonnegative inequality
multipliers, exact right side and all $5996$ column slacks are
independently replayed over the rationals, without an optimizer.

The nine $A_E$ occur once in $I(1)$, since the source ring is
commutative. Adjunction gives
$\sum_E n_E=[F(B)^2:1]=5$. Selfduality of $B$ makes the unit
contribution $1$; the bounds give at most $1+1+1=3$ in total,
a contradiction.
\end{proof}

\subsection{A strict bosonic-row bound}
\begin{lemma}\label{lem:r137-bosonic-mass}
Let $s$ be symmetric unitary normalized modular data with actual
twists $t$ and anomaly $\xi$. Let $J$ be twist-one labels and $X\notin J$,
with $s(X,j)$ real for $j\in J$. Then
$m=\sum_{j\in J}s(X,j)^2\le3/4$.
If $J$ is duality-closed and equality holds, then $X=X^*$,
$\theta_X=-1$, $\xi=\pm1$, and $s(X,X)=\xi/2$.
\end{lemma}
\begin{proof}
The diagonal modular equation and row norm give
$m\le1-m+|s(X,X)|\le1-m+\sqrt{1-m}$.
This is equivalent to $m\le3/4$.
At equality the entire outside-$J$ norm is at $s(X,X)$.
If $X\ne X^*$, the two conjugate rows have the same real restriction
to $J$ of squared norm $3/4$, and distinct remaining diagonal supports,
contradicting orthogonality. Thus $s(X,X)=\pm1/2$, and
$3/4+\theta_X/4=\xi\theta_X^{-2}s(X,X)$ forces the stated values.
\end{proof}

For R137 use all nine distinguished $A_E$ as $J$. A noncanonical
restriction has $v_0=0$. Equality, with center anomaly one and
orthogonality to the unit row, would give
\[
\sum_E\frac{\chi_E(v)}{f_E^2}+\frac{a+bd}{2D}=0,
\quad b=2a,\quad v_2+v_7=3v_1+4v_8.
\]
With $p=v_1,q=v_8$, exact substitution makes its mass
\[
\frac{9p^2+24pq+21q^2+(v_2-v_7)^2+(v_3-v_6)^2+(v_4-v_5)^2}{16}
+\frac{(v_3+v_6-v_4-v_5)^2}{64}.
\]
Equality $3/4$ forces $p=1,q=0$, $(v_2-v_7)^2=1$, and
$v_3+v_4+v_5+v_6=8$. The residual equality would be
\[
(v_3-v_6)^2+(v_4-v_5)^2+
\frac{(v_3+v_6-v_4-v_5)^2}{4}=2.
\]
If $v_3+v_6$ and $v_4+v_5$ are even the left side is $0\bmod4$;
if odd it is $3\bmod4$. Hence the mass is \emph{strictly} less than
$3/4$ for every nonzero noncanonical integral restriction vector.
This proof uses no bound on vector size.

\paragraph{The surviving negative-sign system.}
An exact rational witness of support $124$ satisfies all $256$ general
inherited equations after removing the eight chosen indicator test
equations, full-column Galois pairing, all $45$ nonnegative integral
entries and vanishing irrational parts in the distinguished fusion-by-$B$
block, and the strict mass bound. Its four projected third indicators
are $1/2$ and fourth indicators are $2$; these are witness choices,
not forced values for all negative-sign categories. A separate verifier
rebuilds the field arithmetic and checks every supported equation.
The witness is a solution of necessary conditions, not integral
central data, modular matrices, associators or a categorification.
The complete categorical domain and positive-sign cubic theorem
remain premises of Theorem~\ref{thm:r137-positive-ty-exclusion}.
Current proofs and portable replay are in
\path{evidence/frontier/replay/current/r137/} and the independent
\path{unified/} checks alongside it.
